\subsection{Next combination in Gray code order}\label{sec:successor}

\begin{algorithm}[t!]
\caption{Next combination in Gray order}\label{alg:successor}
\begin{algorithmic}[1]
\Statex{\hspace*{-\algorithmicindent} \textbf{Input:} $a=(a_1,\dots,a_{n})$: $n$-tuple, $n>1$, total sum $d<n$}
\Statex{\hspace*{-\algorithmicindent} \textbf{Output:} $b=(b_1,\dots,b_{n})$ next $n$-tuple combination in Gray order, or $\varnothing$}
% \Statex{\hspace*{-\algorithmicindent} \textbf{Uses:} \texttt{after}, \texttt{call}: , \texttcalt{next}, \texttt{prev}: direction $mode$, one the opposite of the other ($\overline{\texttt{next}}$ denotes \texttt{prev} and vice versa)}
\Statex{\hspace*{-\algorithmicindent} \textbf{Uses:} 
\texttt{call} (simulate recursive call), 
\texttt{after} (action after simulated recursion), 
\texttt{next}/\texttt{prev} (direction to scan the current component; 
$\overline{\texttt{next}}$ denotes \texttt{prev} and vice versa)}

\Function{ComputeTail}{$m$, $s$, parity}
  \State{\IfThen{$m \le 0$}{\Return{$\varnothing$}}}
  \State{\IfThenElse{parity = 0}{\Return{$(\mathbf{0}^{m-1}) || (s)$}}{\Return{$(s) || (\mathbf{0}^{m-1})$}}}
\EndFunction{}
\Function{Successor}{$a$, $n$, $d$}
\State{\IfThen{$a = \varnothing$}{\Return{$(\mathbf{0}^{n-1}, d)$}}}\Comment{First combination in Gray code order}


\State{$P_k=\sum_{j<k} a_j$, for $1\le k\le n$}\Comment{Compute prefix sums}
\State{\texttt{stack} $\gets$ $(\texttt{call},\texttt{next},0,-1)$}\Comment{Initialize stack}
\State{\texttt{last\_ret} $\gets \varnothing$}

\While{\texttt{stack} not empty}

    \State{$(type,mode,k,i)$ $\gets$ \texttt{stack}.pop()}

    \If{$type=\texttt{call}$}
        \If{$k=n-1$}
            \State{\texttt{last\_ret} $\gets \varnothing$}
        \Else{}
            \State{$i \gets a_k$}
            \State{\texttt{stack}.push$((\texttt{after},mode,k,i))$}
            \State{$mode$ $\gets$ \GetsIfThenElse{$mode$}{$i \bmod 2 = 0$}{$\overline{mode}$}}
            \State{\texttt{stack}.push$((\texttt{call},mode,k+1, -1))$ }\Comment{Simulate recursive call at level $k+1$}
        \EndIf{}

    \Else{}
        \State{\texttt{rem} $\gets d-P_k$}

        \If{\texttt{last\_ret} $\neq \varnothing$}
            \State{\texttt{last\_ret} $\gets (i)\ ||\ \texttt{last\_ret}$}
        \Else{}
            \State{\IfThenElse{$mode=\texttt{next}$}{$i \gets i+1$}{$i \gets i-1$}}
            \If{$0 \le i \le \texttt{rem}$}
                \State{\texttt{tail} $\gets$ \Call{ComputeTail}{$n-k-1$, $\texttt{rem}-i$, $\gets i \bmod 2$}}
                \State{\texttt{last\_ret} $\gets (i)\ ||\ \texttt{tail}$}
            \EndIf{}
        \EndIf{}
    \EndIf{}

\EndWhile{}

\State{\Return{\texttt{last\_ret}}}
\EndFunction{}
\end{algorithmic}
\end{algorithm}

% \ma{Old description, change both this and section 4 part, referring to Appendix.}
% As shown in~\cite{DBLP:journals/cj/Takaoka99,DBLP:journals/corr/Takaoka15a}, the Gray code for combinations of a multiset~\cite{DBLP:journals/ejc/RuskeyS96} can be generated efficiently, either via a recursive procedure or its iterative variant. Moreover, these results provide a method to determine the successor of any given combination in constant time. Consequently, the admissible pairs can be generated alongside the corresponding $S$-polynomial computations, without increasing the overall computational complexity of \Cref{alg:step1}.


Our construction follows the general idea underlying Gray code generation for multiset combinations as developed in~\cite{DBLP:journals/ejc/RuskeyS96,DBLP:journals/cj/Takaoka99,DBLP:journals/corr/Takaoka15a}, which we adapt to the requirements of our setting. 
In particular, we define in \Cref{alg:successor} the \textsc{Successor} procedure that, given an input tuple, returns the next tuple in the Gray code ordering, equivalently along the corresponding Hamiltonian path. 
This allows admissible pairs to be generated directly from the traversal of the Gray code ordering while preserving the efficiency of the underlying construction. 
It has been proved that the successor of a combination can be computed in constant time, and our adaptation retains this property~\cite{DBLP:journals/cj/Takaoka99,DBLP:journals/corr/Takaoka15a}.
Since the procedure is obtained as a direct modification of the constructions presented in the cited works, we do not address its correctness here, which follows from the same arguments.

For clarity, \Cref{alg:successor} uses stack labels \texttt{call} and \texttt{after} to simulate the recursive construction, and \texttt{next}/\texttt{prev} to indicate the direction of the component scan.
The procedure internally calls \textsc{ComputeTail} to generate the tail of the tuple (zero-padded entries) according to the remaining sum and parity.  
For further details, we refer the reader to Algorithm 2 and Section 7 in~\cite{DBLP:journals/corr/Takaoka15a}.
